\chapter*{原书序}
\addcontentsline{toc}{chapter}{原书序}
\markboth{原书序}{原书序}

书名 All of Statistics 按字面夸大，按意图则是：在一本入门课里覆盖比传统数理统计更广的现代题目
（Bootstrap、曲线估计、图模型、分类），并尽快把概率接到推断上。
读者预设只会微积分和一点线性代数。

概率研究的正向问题是：给定数据生成过程，结果有什么性质。
统计推断是它的逆问题：给定结果，能对生成过程说什么。
预测、分类、聚类、估计都是后一问题的特例。

\bookfig{fig-preface-probability-inference.png}{1}{概率与推断。}{Data generating process：数据生成过程；Observed data：观测数据；Probability：由过程推结果；Inference and Data Mining：由结果推过程。}

全书三部分：第~I~部分是作为推断语言的概率；
第~II~部分是推断及其近亲数据挖掘、机器学习；
第~III~部分把第~II~部分用到回归、图、因果、密度估计、光滑、分类与模拟，
并补一章随机过程。

统计学与计算机科学常用不同词说同一件事。原书给出对照表如下。

\begin{table}[htbp]
\centering
\small
\begin{tabularx}{\linewidth}{>{\raggedright\arraybackslash}p{3.6cm}>{\raggedright\arraybackslash}p{4.2cm}X}
\toprule
统计学 & 计算机科学 & 含义\\
\midrule
估计（estimation） & 学习（learning） &
用数据估计未知量\\
分类（classification） & 监督学习（supervised learning） &
由 $X$ 预测离散 $Y$\\
聚类（clustering） & 无监督学习（unsupervised learning） &
把数据分成若干组\\
数据（data） & 训练样本（training sample） &
$(X_1,Y_1),\ldots,(X_n,Y_n)$\\
协变量（covariates） & 特征（features） &
各个 $X_i$\\
分类器（classifier） & 假设（hypothesis） &
从协变量到结果的映射\\
假设（hypothesis） & --- &
参数空间 $\Theta$ 的子集\\
置信区间（confidence interval） & --- &
以指定频率覆盖未知量的区间\\
有向无环图（DAG） & 贝叶斯网（Bayes net） &
带指定条件独立关系的多元分布\\
贝叶斯推断 & 贝叶斯推断 &
用数据更新信念的方法\\
频率派推断（frequentist inference） & --- &
有频率保证的统计方法\\
大偏差界（large deviation bounds） & PAC 学习 &
误差概率的一致上界\\
\bottomrule
\end{tabularx}
\end{table}

\noindent 符号：$X$ 为协变量或输入；$Y$ 为结果或标签；
$\Theta$ 为参数空间；DAG 为有向无环图。
表中「含义」列说明该术语在问题中代表什么对象，不是又一个待估参数。
